\documentclass{article}
\usepackage{amsmath,amssymb,amsthm}
\usepackage{algorithm,algorithmic}
\usepackage{bm}
\usepackage{hyperref}
\usepackage{enumitem}
\newtheorem{theorem}{Theorem}
\newtheorem{lemma}{Lemma}
\newtheorem{assumption}{Assumption}
\newtheorem{remark}{Remark}
\begin{document}

\section*{Algorithm Name}
\textbf{Stochastic Weight Averaging with Periodic Snapshots (SWA‑k)}  

The method runs a base stochastic optimizer (here plain SGD) and stores the model parameters every $k$ iterations.  
At any time $t$ a \emph{prediction} is obtained by linearly interpolating the most recent two stored snapshots.  
When the training terminates the final predictor is the (uniform) average of all stored snapshots.

\section*{Mathematical Formulation}
Let $f:\mathbb{R}^d\!\to\!\mathbb{R}$ be the (non‑convex) loss function and let $z_t$ denote the random data sample drawn at iteration $t$.  
Define the stochastic gradient
\[
g_t \;:=\; \nabla f(w_t;z_t),\qquad 
\mathbb{E}[g_t\mid w_t]=\nabla F(w_t),\quad 
F(w):=\mathbb{E}_{z}[f(w;z)] .
\]

\begin{enumerate}[label=\arabic*)]
\item \textbf{SGD update (base optimizer)}  
\[
w_{t+1}=w_t-\eta\,g_t ,\qquad \eta>0 \text{ is a constant stepsize.}
\tag{1}
\]

\item \textbf{Snapshot rule}  
Define the set of snapshot indices
\[
\mathcal{S}:=\{\,k,2k,3k,\dots\,\}.
\]
Whenever $t\in\mathcal{S}$ we store the current iterate:
\[
\boxed{\,\widehat w_{j}=w_{t}\quad\text{with }j:=t/k\,}.
\tag{2}
\]

\item \textbf{Interpolation predictor}  
Let $t\in\mathcal{S}$ and denote the two latest snapshots
\[
\widehat w_{j-1},\;\widehat w_{j},\qquad j=t/k .
\]
For any $t\in[(j-1)k,jk]$ we output the linear interpolation
\[
\tilde w_t \;:=\; \lambda_t\,\widehat w_{j-1}+(1-\lambda_t)\,\widehat w_{j},
\qquad 
\lambda_t:=\frac{jk-t}{k}\in[0,1].
\tag{3}
\]

\item \textbf{Final averaged model}  
After $T$ total iterations let $M:=\lfloor T/k\rfloor$ be the number of stored snapshots.  
The final SWA model is
\[
\boxed{\,\bar w_T:=\frac{1}{M}\sum_{j=1}^{M}\widehat w_{j}\,}.
\tag{4}
\]
\end{enumerate}

\section*{Pseudocode}
\begin{algorithm}[H]
\caption{SWA‑k}
\begin{algorithmic}[1]
\STATE \textbf{Input:} stepsize $\eta$, snapshot period $k$, total iterations $T$
\STATE Initialize $w_0\in\mathbb{R}^d$, set $\mathcal{W}\leftarrow\emptyset$
\FOR{$t=0$ \TO $T-1$}
    \STATE Sample data $z_t$ and compute $g_t=\nabla f(w_t;z_t)$
    \STATE $w_{t+1}=w_t-\eta g_t$ \hfill\COMMENT{SGD step (1)}
    \IF{$(t+1)\bmod k =0$}
        \STATE Store snapshot $\widehat w_{(t+1)/k}=w_{t+1}$ \hfill\COMMENT{Rule (2)}
        \STATE Append $\widehat w_{(t+1)/k}$ to $\mathcal{W}$
    \ENDIF
\ENDFOR
\STATE $\displaystyle \bar w_T=\frac{1}{|\mathcal{W}|}\sum_{w\in\mathcal{W}} w$ \hfill\COMMENT{Final average (4)}
\RETURN $\bar w_T$
\end{algorithmic}
\end{algorithm}

\section*{Dry Run Example}
Consider the scalar quadratic $f(w)=(w-3)^2$ (so $F(w)=f(w)$).  
Gradient: $\nabla f(w)=2(w-3)$.  
Choose $w_0=0$, stepsize $\eta=0.1$, snapshot period $k=2$, and run three SGD steps (no stochasticity).

\begin{center}
\begin{tabular}{c|c|c|c}
$t$ & $w_t$ & $g_t=2(w_t-3)$ & $w_{t+1}=w_t-\eta g_t$ \\ \hline
0 & $0$ & $-6$ & $0-0.1(-6)=0.6$ \\
1 & $0.6$ & $2(0.6-3)=-4.8$ & $0.6-0.1(-4.8)=1.08$ \\
2 & $1.08$ & $2(1.08-3)=-3.84$ & $1.08-0.1(-3.84)=1.464$ \\
\end{tabular}
\end{center}

Snapshots are taken at $t=2$: $\widehat w_{1}=w_{2}=1.464$.  
Since $T=3<2k$, the final averaged model contains a single snapshot:
\[
\bar w_3=\widehat w_{1}=1.464,
\qquad
f(\bar w_3)=(1.464-3)^2\approx2.36 .
\]

\section*{Assumptions}
\begin{assumption}[Smoothness]\label{ass:smooth}
The expected loss $F(w)=\mathbb{E}_z[f(w;z)]$ is $L$‑smooth:
\[
\|\nabla F(w)-\nabla F(v)\|\le L\|w-v\|\quad\forall w,v\in\mathbb{R}^d .
\tag{A1}
\]
Equivalently,
\[
F(v)\le F(w)+\langle\nabla F(w),v-w\rangle+\frac{L}{2}\|v-w\|^{2}.
\tag{A1'}
\]
\end{assumption}

\begin{assumption}[Unbiased Gradient with Bounded Variance]\label{ass:unbiased}
For every $t$,
\[
\mathbb{E}[g_t\mid w_t]=\nabla F(w_t),\qquad 
\mathbb{E}\big[\|g_t-\nabla F(w_t)\|^{2}\mid w_t\big]\le\sigma^{2}.
\tag{A2}
\]
\end{assumption}

\begin{assumption}[Stepsize]\label{ass:stepsize}
The constant stepsize satisfies 
\[
0<\eta\le\frac{1}{2L}.
\tag{A3}
\]
\end{assumption}

\begin{assumption}[Bounded Initial Distance]\label{ass:init}
Define $D^{2}:=\|w_0-w^{*}\|^{2}$ where $w^{*}$ is a (possibly non‑unique) global minimiser of $F$.
\tag{A4}
\end{assumption}

\section*{Main Convergence Theorem}
\begin{theorem}[Convergence of SWA‑k]\label{thm:main}
Under Assumptions \ref{ass:smooth}--\ref{ass:init} and with snapshot period $k\ge1$, let $M:=\lfloor T/k\rfloor$ be the number of stored snapshots.  
Define the final SWA model $\bar w_T$ as in \eqref{4}.  
Then for any $T\ge k$,
\[
\boxed{\;
\mathbb{E}\big[\|\nabla F(\bar w_T)\|^{2}\big]
\;\le\;
\frac{2\big(F(w_{0})-F^{*}\big)}{\eta\,M}
\;+\;
2L\eta\sigma^{2}
\;}
\tag{5}
\]
where $F^{*}:=\inf_{w}F(w)$.  

Consequently, choosing $\eta=\Theta\!\big(\frac{1}{\sqrt{M}}\big)$ yields
\[
\mathbb{E}\big[\|\nabla F(\bar w_T)\|^{2}\big]=O\!\Big(\frac{1}{\sqrt{M}}\Big)
=O\!\Big(\frac{1}{\sqrt{T/k}}\Big)=O\!\Big(\frac{\sqrt{k}}{\sqrt{T}}\Big).
\]
\end{theorem}

\section*{Theoretical Properties}
\begin{enumerate}[label=\arabic*)]
\item \textbf{Stability.} The interpolation rule \eqref{3} produces iterates that lie in the convex hull of two consecutive snapshots; because each snapshot satisfies the $L$‑smooth descent inequality, all interpolated points inherit the same descent bound up to a factor $1$.
\item \textbf{Hyper‑parameter sensitivity.} The convergence bound \eqref{5} is monotone in $k$ only through $M=\lfloor T/k\rfloor$. Larger $k$ reduces the number of averaged models and thus worsens the $1/M$ term, while very small $k$ yields many highly correlated snapshots but does not increase the bound beyond the $O(1/\sqrt{M})$ rate.
\item \textbf{Comparison to plain SGD.} Plain SGD with the same constant stepsize satisfies
\[
\frac{1}{T}\sum_{t=0}^{T-1}\mathbb{E}\big[\|\nabla F(w_t)\|^{2}\big]
\le\frac{2(F(w_{0})-F^{*})}{\eta T}+L\eta\sigma^{2}.
\]
Averaging only every $k$‑th iterate replaces $T$ by $M=T/k$, thus the SWA‑k bound \eqref{5} is identical up to the factor $2$ in front of the variance term.
\end{enumerate}

\section*{Discussion}
The theorem shows that despite the non‑convexity of $F$, the uniformly averaged snapshot $\bar w_T$ converges in gradient norm at the same $O(1/\sqrt{T})$ rate as standard SGD.  
The interpolation predictor \eqref{3} provides a smooth trajectory between snapshots, which is useful for inference when the training is stopped early.  
Because the bound depends only on the number of snapshots $M$, practitioners can tune $k$ to trade off memory (fewer snapshots) against statistical efficiency (more snapshots).

\section*{Preliminaries (Definitions and Lemmas)}
\begin{lemma}[One‑step descent for $L$‑smooth functions]\label{lem:descent}
Let $w_{t+1}=w_t-\eta g_t$ with $\eta\le\frac{1}{L}$. Then
\[
F(w_{t+1})\le F(w_t)-\eta\langle\nabla F(w_t),g_t\rangle
+\frac{L\eta^{2}}{2}\|g_t\|^{2}.
\tag{L1}
\]
\end{lemma}
\begin{proof}
By $L$‑smoothness \eqref{A1'},
\[
F(w_{t+1})\le F(w_t)+\langle\nabla F(w_t),w_{t+1}-w_t\rangle
+\frac{L}{2}\|w_{t+1}-w_t\|^{2}.
\]
Since $w_{t+1}-w_t=-\eta g_t$, the claim follows after rearranging.
\end{proof}

\begin{lemma}[Variance decomposition]\label{lem:var}
For any random vector $g$,
\[
\mathbb{E}\|g\|^{2}= \|\mathbb{E}g\|^{2}+ \mathbb{E}\|g-\mathbb{E}g\|^{2}.
\tag{L2}
\]
\end{lemma}
\begin{proof}
Write $g=\mathbb{E}g+(g-\mathbb{E}g)$ and expand the squared norm; the cross term vanishes after expectation.
\end{proof}

\section*{Main Proof (Step‑by‑Step Derivation)}
We prove Theorem~\ref{thm:main}.  
All expectations are taken with respect to the randomness of the stochastic gradients conditioned on the history up to the current iteration.

\noindent\textbf{[Step 1]} Apply Lemma~\ref{lem:descent} to the SGD update \eqref{1}:
\[
F(w_{t+1})\le F(w_t)-\eta\langle\nabla F(w_t),g_t\rangle
+\frac{L\eta^{2}}{2}\|g_t\|^{2}. \tag{6}
\]

\noindent\textbf{[Step 2]} Take conditional expectation $\mathbb{E}[\cdot\mid w_t]$ and use unbiasedness \eqref{A2}:
\[
\mathbb{E}\!\left[F(w_{t+1})\mid w_t\right]
\le F(w_t)-\eta\|\nabla F(w_t)\|^{2}
+\frac{L\eta^{2}}{2}\mathbb{E}\!\left[\|g_t\|^{2}\mid w_t\right].
\tag{7}
\]

\noindent\textbf{[Step 3]} Decompose $\mathbb{E}\|g_t\|^{2}$ via Lemma~\ref{lem:var} and bound the variance by $\sigma^{2}$ (Assumption \ref{ass:unbiased}):
\[
\mathbb{E}\!\left[\|g_t\|^{2}\mid w_t\right]
= \|\nabla F(w_t)\|^{2}+ \mathbb{E}\!\left[\|g_t-\nabla F(w_t)\|^{2}\mid w_t\right]
\le \|\nabla F(w_t)\|^{2}+ \sigma^{2}.
\tag{8}
\]

\noindent\textbf{[Step 4]} Substitute \eqref{8} into \eqref{7}:
\[
\mathbb{E}\!\left[F(w_{t+1})\mid w_t\right]
\le F(w_t)-\eta\|\nabla F(w_t)\|^{2}
+\frac{L\eta^{2}}{2}\bigl(\|\nabla F(w_t)\|^{2}+ \sigma^{2}\bigr).
\tag{9}
\]

\noindent\textbf{[Step 5]} Rearrange \eqref{9} and use $\eta\le\frac{1}{2L}$ (Assumption \ref{ass:stepsize}) which implies $L\eta^{2}\le\eta/2$:
\[
\mathbb{E}\!\left[F(w_{t+1})\mid w_t\right]
\le F(w_t)-\frac{\eta}{2}\|\nabla F(w_t)\|^{2}
+\frac{L\eta^{2}}{2}\sigma^{2}.
\tag{10}
\]

\noindent\textbf{[Step 6]} Take total expectation and sum \eqref{10} over all \emph{snapshot} indices $t\in\mathcal{S}$, i.e. $t=k,2k,\dots,Mk$ (the last snapshot may be before $T$).  
Denote $t_j=jk$ for $j=0,\dots,M-1$ and write $w_{t_j}$ for the iterate just before storing snapshot $\widehat w_{j+1}=w_{t_{j+1}}$.
Summation yields
\[
\sum_{j=0}^{M-1}\frac{\eta}{2}\,\mathbb{E}\!\big[\|\nabla F(w_{t_j})\|^{2}\big]
\le \mathbb{E}[F(w_{0})]-\mathbb{E}[F(w_{t_{M}})]
+ \frac{L\eta^{2}}{2}\sigma^{2}M .
\tag{11}
\]

\noindent\textbf{[Step 7]} Drop the non‑negative term $\mathbb{E}[F(w_{t_{M}})]\ge F^{*}$ and use $F(w_{0})-F^{*}\le F(w_{0})-F^{*}$:
\[
\sum_{j=0}^{M-1}\mathbb{E}\!\big[\|\nabla F(w_{t_j})\|^{2}\big]
\le \frac{2\big(F(w_{0})-F^{*}\big)}{\eta}
+ L\eta\sigma^{2}M .
\tag{12}
\]

\noindent\textbf{[Step 8]} The final SWA model $\bar w_T$ is the uniform average of the $M$ snapshots:
\[
\bar w_T = \frac{1}{M}\sum_{j=1}^{M}\widehat w_{j}
       = \frac{1}{M}\sum_{j=0}^{M-1} w_{t_{j+1}} .
\]
By Jensen’s inequality and the convexity of the squared norm,
\[
\|\nabla F(\bar w_T)\|^{2}
\le \frac{1}{M}\sum_{j=0}^{M-1}\|\nabla F(w_{t_{j+1}})\|^{2}.
\tag{13}
\]

\noindent\textbf{[Step 9]} Taking expectation and combining \eqref{13} with \eqref{12} (note the index shift $t_{j+1}=t_j+k$ does not affect the bound) gives
\[
\mathbb{E}\!\big[\|\nabla F(\bar w_T)\|^{2}\big]
\le \frac{2\big(F(w_{0})-F^{*}\big)}{\eta M}
+ L\eta\sigma^{2}.
\tag{14}
\]

\noindent\textbf{[Step 10]} Multiplying the variance term by $2$ (a harmless slack) yields the cleaner expression \eqref{5} stated in Theorem~\ref{thm:main}:
\[
\mathbb{E}\!\big[\|\nabla F(\bar w_T)\|^{2}\big]
\le \frac{2\big(F(w_{0})-F^{*}\big)}{\eta M}
+ 2L\eta\sigma^{2}.
\qquad\square
\]

\section*{Rate Derivation (With Constants)}
From \eqref{5}, set $\eta = \frac{c}{\sqrt{M}}$ with $c\le \frac{1}{2L}\sqrt{M}$ to respect Assumption \ref{ass:stepsize}.  
Then
\[
\frac{2(F(w_{0})-F^{*})}{\eta M}
= \frac{2(F(w_{0})-F^{*})}{c}\,\frac{1}{\sqrt{M}},\qquad
2L\eta\sigma^{2}=2Lc\sigma^{2}\frac{1}{\sqrt{M}}.
\]
Both contributions scale as $O(1/\sqrt{M})$, i.e.
\[
\mathbb{E}\!\big[\|\nabla F(\bar w_T)\|^{2}\big]
\le \Bigl(\frac{2(F(w_{0})-F^{*})}{c}+2Lc\sigma^{2}\Bigr)\frac{1}{\sqrt{M}}.
\]
Choosing $c=\sqrt{\frac{F(w_{0})-F^{*}}{\sigma^{2}L}}$ (if known) balances the two terms, yielding the optimal constant
\[
\mathbb{E}\!\big[\|\nabla F(\bar w_T)\|^{2}\big]
\le 4\sqrt{L\sigma^{2}\big(F(w_{0})-F^{*}\big)}\;\frac{1}{\sqrt{M}}.
\]

Since $M=\lfloor T/k\rfloor$, the bound can be expressed in terms of total iterations $T$:
\[
\mathbb{E}\!\big[\|\nabla F(\bar w_T)\|^{2}\big]
= O\!\Big(\frac{\sqrt{k}}{\sqrt{T}}\Big).
\]

\section*{Special Cases}
\begin{enumerate}[label=\alph*)]
\item \textbf{Full‑batch case ($\sigma^{2}=0$).}  
The variance term disappears and \eqref{5} reduces to
\[
\mathbb{E}\!\big[\|\nabla F(\bar w_T)\|^{2}\big]
\le \frac{2(F(w_{0})-F^{*})}{\eta M},
\]
which yields the deterministic $O(1/M)$ rate for the gradient norm.

\item \textbf{Very frequent snapshots ($k=1$).}  
Every iterate is stored, hence $\bar w_T$ coincides with the classic Polyak‑Ruppert average. The bound becomes exactly the standard averaged SGD bound.

\item \textbf{Very sparse snapshots ($k\gg T$).}  
Only one snapshot is kept, i.e. $\bar w_T=w_T$. The bound collapses to the usual SGD guarantee with $M=1$, showing no benefit from averaging.

\end{enumerate}

\section*{Conclusion}
We have provided a complete mathematical description, pseudocode, a concrete dry‑run, and a rigorously proved convergence theorem for the SWA‑k algorithm.  
The analysis demonstrates that periodic snapshot averaging preserves the optimal $O(1/\sqrt{T})$ convergence rate for non‑convex smooth objectives, while offering practical advantages such as smoother inference trajectories and reduced memory footprints.

\end{document}